\section{Phương pháp và kiến trúc hệ thống}\label{sec:method}

\subsection{Phương pháp nghiên cứu}
Nghiên cứu áp dụng phương pháp khoa học thiết kế (\emph{Design Science Research}, DSR)~\parencite{hevner2004design}. Theo DSR, tri thức được tạo ra thông qua việc xây dựng và đánh giá một hiện vật kỹ thuật có chủ đích. Hiện vật của đề tài là toàn bộ kiến trúc Continux: cụm K3s, luồng dữ liệu từ Vector đến Iceberg, lớp quản lý triển khai Argo CD và lớp quan sát VictoriaMetrics/Grafana.

Đánh giá được thực hiện bằng bốn lượt chạy độc lập trên cùng bản sửa đổi Git \texttt{ff19448}. Mỗi lượt đi qua các pha: khởi tạo, chuẩn bị dữ liệu, tạo trạng thái nền sạch, phát lại, chuyển giao MV, xác minh cuối, dọn trạng thái chạy và dọn tệp cục bộ. Cách tổ chức này giúp mỗi cấu hình bắt đầu từ cùng một trạng thái và tránh trộn dữ liệu giữa các lượt.

\subsection{Dữ liệu và cấu hình phát lại}
Sở Taxi và Dịch vụ vận tải New York (NYC TLC) công bố dữ liệu chuyến xe công khai phục vụ phân tích~\parencite{nyctlc}. Continux dùng Yellow Taxi tháng \texttt{2026-03}. Tập lệnh \path{scripts/partojsonl.py} chuyển Parquet thành JSONL để Vector đọc tuần tự. Bảng Taxi Zone Lookup được tải lên MinIO để ánh xạ mã vị trí sang quận và khu vực.

\begin{table*}[t!]
\centering
\caption{Nguồn dữ liệu và cấu hình phát lại}
\label{tab:dataset-profiles-vi}
\begin{tabularx}{\textwidth}{@{}p{0.37\textwidth}Y@{}}
\toprule
\textbf{Thuộc tính} & \textbf{Giá trị}\\
\midrule
Dữ liệu chuyến xe & NYC TLC Yellow Taxi, tháng \texttt{2026-03}\\
Định dạng nguồn / đích chuyển đổi & Parquet / JSONL\\
Số dòng JSONL xác nhận & \texttt{3.952.451}\\
Bảng tra cứu Taxi Zone Lookup & \texttt{265} dòng\\
Giới hạn phát cấu hình & \texttt{2}, \texttt{1.000}, \texttt{5.000}, \texttt{10.000} sự kiện/s\\
Thời gian lấy mẫu khi phát lại & \texttt{12} mẫu, cách nhau \texttt{5 s}\\
\bottomrule
\end{tabularx}
\end{table*}

Vector bổ sung \texttt{event\_id} và \texttt{event\_time}, rồi gửi từng sự kiện vào topic \texttt{nyc-taxi-events}. Tham số tốc độ chỉ giới hạn nguồn phát. Vì báo cáo không có phép đo số sự kiện đầu cuối trên mỗi giây, các cấu hình được dùng để so sánh xu hướng số chuyến đã xử lý trong cùng cửa sổ quan sát, không được diễn giải thành thông lượng thực tế.

\subsection{Kiến trúc tổng thể}
\begin{figure*}[t!]
\centering
\includegraphics[width=\textwidth,height=0.72\textheight,keepaspectratio]{architecture-process-flow.png}
\caption{Kiến trúc, luồng xử lý và quy trình chuyển giao của Continux 2.1.0.}
\label{fig:architecture-vi}
\end{figure*}

Luồng dữ liệu bắt đầu từ JSONL, đi qua Vector và Redpanda, sau đó được RisingWave xử lý bằng SQL luồng. RisingWave join sự kiện với bảng tra cứu, duy trì MV và ghi kết quả Iceberg lên MinIO. Luồng điều khiển dùng Argo CD để đồng bộ manifest từ Git. Luồng quan sát dùng exporter, VictoriaMetrics và Grafana để hiển thị trạng thái phát lại, tài nguyên, chuyển giao và tính toàn vẹn quan sát được.

\subsection{Cấu trúc triển khai}
\begin{table*}[t!]
\centering
\caption{Cấu trúc cụm K3s thực nghiệm}
\label{tab:topology-vi}
\begin{tabularx}{\textwidth}{@{}p{0.18\textwidth}p{0.40\textwidth}Y@{}}
\toprule
\textbf{Node} & \textbf{Tài nguyên ghi nhận} & \textbf{Vai trò}\\
\midrule
\texttt{imac} & Intel i5-8500, 6 core, 8 GB RAM, 200 GB SSD, Ubuntu 26.04 & Máy chủ \#1; Vector, Redpanda, RisingWave, MinIO, exporter\\
\texttt{continux-vps} & 2 vCPU, 4 GB RAM, 80 GB SSD, Ubuntu 24.04 & Máy chủ \#2; Argo CD, VictoriaMetrics, Grafana\\
\texttt{helios-pc} & WSL Ubuntu 26.04, host Intel i5-12500H, 16 GB RAM & Máy chủ \#3; quorum-only, node-exporter\\
\bottomrule
\end{tabularx}
\end{table*}

Ba node kết nối qua Tailscale. K3s dùng embedded etcd với quorum \texttt{2/3}. Thuật ngữ sẵn sàng cao trong báo cáo chỉ áp dụng cho lớp điều khiển Kubernetes và etcd. Đường dữ liệu không phải cấu hình sẵn sàng cao đầu cuối: Redpanda chạy một broker và topic dùng \texttt{replicationFactor=1}.

\clearpage
\onecolumn
\subsection{Đối tượng SQL và chuyển giao Blue/Green}
\begin{table}[H]
\centering
\caption{Đối tượng dữ liệu cốt lõi}
\begin{tabularx}{\textwidth}{@{}p{0.27\textwidth}p{0.27\textwidth}Y@{}}
\toprule
\textbf{Đối tượng} & \textbf{Đầu vào} & \textbf{Chức năng}\\
\midrule
\path{tlc_zone} & CSV trên MinIO & Tra cứu vị trí sang quận và khu vực\\
\path{nyc_taxi_src} & Topic Redpanda & Nguồn sự kiện JSON\\
\path{mv_zone_stats_blue} & Nguồn + bảng tra cứu & Cách tính nền Blue, giữ toàn bộ bản ghi hợp lệ ở trạng thái ban đầu\\
\path{mv_zone_stats} & MV Blue ban đầu & Tên công khai ổn định để phía truy vấn luôn gọi cùng một đối tượng\\
\path{mv_zone_stats_green} & Nguồn + tra cứu + bộ lọc & Ứng viên Green, dùng để kiểm tra logic mới trước khi chuyển giao\\
\path{sink_zone_stats} & MV công khai lúc tạo đích ghi & Đầu ra Iceberg kiểu cập nhật hoặc chèn\\
\bottomrule
\end{tabularx}
\end{table}

Cách tính Green bổ sung điều kiện loại các bản ghi có giá trị âm ở tiền cước hoặc quãng đường:

\begin{lstlisting}[language=SQL,caption={Điều kiện bổ sung trong cách tính Green}]
WHERE t.fare_amount >= 0
  AND t.trip_distance >= 0
\end{lstlisting}

Khi MV Green đã có dữ liệu ổn định, vòng lặp truy vấn tiếp tục đọc tên công khai theo chu kỳ khoảng \texttt{0,5 s}. Trong lúc vòng lặp này đang chạy, hệ thống thực hiện câu lệnh:

\begin{lstlisting}[language=SQL,caption={Câu lệnh hoán đổi tên MV}]
ALTER MATERIALIZED VIEW mv_zone_stats
SWAP WITH mv_zone_stats_green;
\end{lstlisting}

Phép hoán đổi giữ nguyên tên truy vấn công khai nhưng thay cách tính phía sau tên đó. Nhờ vậy, người dùng vẫn truy vấn cùng một tên MV, còn hệ thống chuyển từ logic Blue sang logic Green. Đích ghi Iceberg đã được tạo trước thao tác này; vì vậy đầu ra Iceberg chứng minh đường ghi hoạt động trong pha phát lại, không chứng minh đích ghi tự chuyển sang cách tính Green.

\subsection{Nguồn bằng chứng và khả năng truy vết}
Bốn gói bằng chứng chính thức được sao chép từ iMac về Windows tại \path{evidence/finalize/20260531-four-profiles/}. Thư mục \path{runs/} chứa bốn \texttt{RUN\_ID}: \texttt{20260531-172050}, \texttt{20260531-173434}, \texttt{20260531-174621} và \texttt{20260531-175701}. Mỗi gói có \texttt{41} tệp; tổng cộng có \texttt{164} tệp văn bản hoặc JSON.

Tệp \path{SHA256SUMS.txt} chứa giá trị băm SHA-256 của từng tệp. Giá trị băm nguồn trên iMac và bản sao Windows khớp \texttt{164/164}. Tệp \path{summary.csv} tổng hợp số liệu chính; tệp \path{sensitive-scan.txt} ghi kết quả rà chuỗi nhạy cảm. Các nhóm bằng chứng được dùng như sau:

\begin{table}[H]
\centering
\caption{Nhóm tệp bằng chứng và mục đích đối chiếu}
\begin{tabularx}{\textwidth}{@{}p{0.53\textwidth}Y@{}}
\toprule
\textbf{Nhóm} & \textbf{Nội dung đối chiếu}\\
\midrule
\path{00-run-id.txt} & Cấu hình và giới hạn phát cấu hình\\
\path{04-mv-progress.txt} & Đường tăng số vùng và số chuyến trong pha phát lại\\
\path{04-mv-final.txt} & Số vùng và số chuyến chốt sau phát lại\\
\path{04-iceberg-final.txt} & Đối tượng Parquet và siêu dữ liệu Iceberg\\
\path{05-green-ready-samples.txt} & Trạng thái MV Green trước chuyển giao\\
\path{05-query-loop-during-cutover.txt} & Truy vấn thành công trước và sau hoán đổi\\
\path{05-duration-and-errors.txt} & Thời lượng hoán đổi và số lỗi truy vấn\\
\path{05-final-*} & SQL cuối, Vector, sức khỏe cụm và Argo CD\\
\path{06-*} & Kết quả dọn trạng thái sinh ra\\
\bottomrule
\end{tabularx}
\end{table}
